\section{The contraction interface and encoding conventions}
\label{sec:interface}

\begin{hypothesis}[\HTC: OpenAI, Theorem D.1]
Fix positive constants $C_0,A_0,b_0$. For an integer $s\ge2$, let a
bipartite multigraph carry one Boolean variable on each edge. There
are at most $C_0s$ edges and a number of vertices bounded by a fixed
polynomial in $s$, on which the constants below may also depend, and its
maximum degree is at most $\lceil s^{9/10}\rceil$. If the total graph
and tensor description length is $\mathcal L$, the tensor entries lie
in
\[
 \mathcal B_{P_0,\ell}=(\Z/2^{P_0}\Z)[V]/(V^\ell),
 \qquad 1\le P_0,\ell\le C_0s.
\]
Uniformly for every vertex and every Boolean port assignment, an entry
can be computed from its description in time
$(s+\mathcal L)^{b_0}2^{A_0s^{9/10}}$ and space
$(s+\mathcal L)^{b_0}$. Then the contraction
\begin{equation}\label{eq:contraction}
 Z=\sum_{x\in\{0,1\}^{E}}\prod_v T_v(x|_{\delta(v)})
\end{equation}
is computable deterministically in
$(s+\mathcal L)^{O(1)}2^{O(g(s))}$ bit operations and
$(s+\mathcal L)^{O(1)}$ space. Constants depend only on the fixed
bounds just specified. Isolated vertices contribute their scalar
entries; $\ell=1$ is allowed.
\end{hypothesis}

This is Theorem~D.1 (Boolean tensor contraction) of~\cite{OpenAI26}.
All our applications use $\ell=1$. The only
coefficients required are integers modulo a power of two, including
negative residues for determinants.

For comparison, write \HBB\ for Theorem 1.1 of~\cite{OpenAI26}:
for each fixed $k\ge2$, exact counting for $k$-CNF on $n\ge20$
declared variables and $m$ clauses takes polynomial space and time
\[
 (n+m)^{O_k(1)}
 2^{O_k(n\log\log n/\sqrt{\log n})}.
\]
For each fixed $D$, if $m\le Dn$, the exponent improves to
$O_{k,D}(g(n))$. The quantifier fixes $k$ and $D$ before the input
size grows.

The accompanying Lean scope document~\cite{OpenAILean26},
\ancfile{lean/docs/378.md}, describes exact counting for fixed $k$ on
canonical encodings, in time
\[
 (L+1)^{O_k(1)}
 2^{O_k(n(\log\log n)^2/\sqrt{\log n})}
\]
and with polynomially many visited tape cells. Its stated scope is
this fixed-$k$ counting assertion. It does not include \HTC, nor the
sharper arbitrary-density or fixed-density bounds in the October 25
Theorem 1.1. Thus that formalization does not certify the hypothesis
used by the results of this paper.

\subsection{Boolean wires and unique extensions}

A functional gadget has input and output ports and internal Boolean
wires. For each assignment to its input ports, exactly one assignment
to its internal wires and output ports has nonzero product of local
$0/1$ entries. A circuit of such gadgets has a unique extension of each
original input, by induction in a topological order. Pinning the final
output to one then makes its contraction the accepting-input count.

Fan-out is handled by equality trees. For a signal with $d\ge2$
occurrences, join its ports by a tree of binary and ternary equality
vertices; fixing the common signal determines every internal bit.
For $d=1$, attach the free unary tensor $(1,1)$. A declared input
with no occurrence is represented by a wire between two copies of
$(1,1)$, contributing exactly two. In particular, changing a gate to a
constant never deletes the factor contributed by an input that thereby
becomes unused. These rules also apply when the circuit output is an
input or a constant.

A constant gate is evaluated directly and its outgoing uses are fixed
by unary pins. A nonconstant gate has at least one input wire, so the
number of remaining gates is at most $w$. Constants, dead gates, and
zero-arity scalar factors can be removed from the network description
after accounting for their effect and retaining all declared inputs.

Given any finite factor graph, insert a binary equality tensor on
each edge. The old and new vertices form a bipartition, edge count
doubles, maximum degree does not increase, and the assignment on each
old edge has exactly one extension to its two replacements. If the
sum of port occurrences is $O(S)$, equality trees and subdivision
therefore use $O(S)$ edges and vertices.
Distinct occurrences at the same factor are distinct ports. If a
two-occurrence signal joins two ports of that factor, subdivision
creates parallel edges to the equality vertex. These are allowed by
the bipartite multigraph interface of \HTC.

\subsection{Scaling and integer readback}

Suppose an encoding has $C_1S$ edges, $O(S)$ vertices, degree at most
$D_1S^{9/10}$, precision $P_0\le C_2S$, and total description length
$\mathcal L$. Suppose its local entry programs use fixed polynomial
time and space in $S+\mathcal L$. Set
\begin{equation}\label{eq:rescale}
 s=\lceil KS\rceil,\qquad
 K\ge\max\{1,(2D_1)^{10/9}\},
\end{equation}
and fix a sufficiently large interface constant $C_0$. All the bounds
in \HTC\ now hold, including its precise maximum-degree condition.
Since $K$ is fixed, $g(s)=O(g(S))$. The same reasoning applies to
$O(S^{2/3})$ degrees. The resulting time is
$(S+\mathcal L)^{O(1)}2^{O(g(S))}$ and the space is polynomial in
$S+\mathcal L$, allowing arbitrary description length. Bounded small
instances may be enumerated with polynomial cost in their descriptions.

If a contraction is known, as an integer, to satisfy
$0\le Z<2^{P_0}$, its least nonnegative residue modulo $2^{P_0}$
recovers $Z$ exactly. For a circuit with unique extensions, $P_0=n+1$
is sufficient, since $0\le Z\le2^n$. For a signed graph encoding the
bound is proved after cancellation, before choosing $P_0$. Division by
a known multiplicity is performed on this recovered integer.
